\section{Introduction}
\label{sec:introduction}

The study of cause and effect has served humanity for millennia and provided a foundation for scientific inquiry and understanding. Contemporary frameworks for computational causality, from Pearl's Structural Causal Models to Granger's time-series analysis, provide a formal basis for causal inference within systems governed by fixed causal structures. These classical models, however, are predicated on a set of core assumptions, including a fixed background spacetime, linear temporal progression, and static causal structures. However, at the frontiers of science and engineering, a new category of challenges arises where the rules of causality themselves become dynamic. For these complex dynamic systems, the classical assumption of a static causal structure embedded in a fixed background spacetime is no longer applicable and imposes a fundamental limitation. From the dynamic regime shifts in financial markets with complex, multi-scale temporal feedback loops to context-dependent safety of autonomous vehicles, there is a need for a new foundation to model systems where causal structures themselves can evolve dynamically.

The presented monograph introduces the Effect Propagation Process (EPP), a single axiomatic foundation for dynamic causal models. Its purpose is to serve as a foundation upon which a new generation of domain-specific dynamic causal models can be built. Categorically, the EPP is closest to a meta-calculus because it defines the abstract mechanisms of dynamic causality while leaving the specifics to a derived dynamic causal model. However, the EPP also borrows from other categories:

\begin{itemize}
\item A Meta-Theory: The EPP provides the philosophical and logical foundation via its metaphysics and ontology upon which one can build a theory of dynamic causality.
\item A Meta-Algebra: The EPP provides the formal, abstract language via its formalization of Causaloids, Contextoids, and their structural relationships.
\item A Meta-Calculus: The EPP provides the formal, computational elements of dynamic causality via the Effect Propagation Process and the Deontic Inference Cycle.
\end{itemize}
  
The EPP adds dynamics as a first class principle to causality based on a single axiomatic foundation that generalizes causality as a spacetime-agnostic functional dependency. The operationalization of the EPP's single axiomatic foundation necessitates a trifecta of computable, first-class primitives: The Context, the Causaloid, and the Causal State Machine.

 The detachment from a fixed background spacetime requires an explicit and dynamic Context. The EPP design of the context enables Euclidean and non-Euclidean representation and linear and non-linear temporal structures, thus supporting the modeling of rich, dynamic, and complex operational environments.

The functional nature of the axiom requires a polymorphic container for any specific causal calculus, the Causaloid. The causaloid, a concept borrowed from physicist Lucian Hardy, unifies cause and effect into a single abstract entity that solves a fundamental problem of structural composition by enabling isomorphic recursive causal structures. The EPP introduces three modalities of dynamic causality: dynamic, adaptive, and emergent. While dynamic and adaptive causality remain deterministic, the introduction of emergent causality, where causal structures co-emerge with their context, also introduces non-determinism with the implication that static verifiability is no longer possible. The Effect Ethos provides an operational guardrail to emergent causality based on a defeasible deontic calculus.

The Causal State Machine is a formal mechanism that translates causal reasoning into actions. A causal state machine separates its state from an action to allow an optional Effect Ethos to verify a proposed action against a set of norms to decide the permissibility of an action relative to the encoded ethos. The causal state machine combined with an Effect Ethos enable dynamic causal system in a dynamic environment while ensuring compliance as code to enforce alignment of derived actions.

\subsection{Precedent}

The presented Effect Propagation Process builds upon rich and diverse preceding scholarly work. 

\textbf{Whitehead's Process Philosophy}

The EPP's primary departure point is a fundamental rejection of the classical Newtonian conception of a static, absolute background spacetime. This move is deeply rooted in the tradition of process philosophy, which argues that reality is not composed of enduring, static substances but is a dynamic flow of interconnected events. This idea finds its clearest expression in the work of Alfred North Whitehead, who posited a universe of "actual occasions"\cite{whitehead2010process}, and Henri Bergson, who described reality as a continuous "creative evolution"\cite{bergson2022creative}. Their shared insight of reality as a process inspires the EPP's foundational redefinition of causality itself, shifting from a static, happen-before relation to a dynamic process of effect propagation.

\newpage

\textbf{Moggi and Wadler's Monadic Composition}
 
The CausalMonad and its bind operation inherit a programming discipline whose theoretical move belongs to Eugenio Moggi and whose practical articulation belongs to Philip Wadler. Moggi proposed in 1989 that monads from category theory provide a useful structuring tool for the denotational semantics of programming languages\cite{moggi1989computational}. His key move was distinguishing a value type $a$ from a computation type $M\ a$ and showing that monads can encapsulate features such as state, exception handling, and continuations within a single uniform interface. Wadler then demonstrated in 1992 that this theoretical insight transfers directly into a programming discipline\cite{wadler1992essence}. A function of type $a \to b$ can be lifted into monadic form $a \to M\ b$, and the resulting programs can be modified to add error handling, state, output, or non-deterministic choice by changing the monad while leaving the program structure essentially intact. The EPP adopts this discipline for causal composition. The monadic axiom $m_2 = m_1 \mathbin{>\!\!>\!\!=} f$ is the Kleisli composition of a monad whose context parameters carry the state, configuration, error condition, and audit log of the causal process. The flexibility that Wadler showed for lifting plain functions into a context-aware monadic form is the same flexibility that allows the CausalMonad to interoperate with the Causaloid through the shared PropagatingEffect type.
 

\textbf{Pearl's SCM}

Judea Pearl, with his SCM, established the foundation upon which the entire field of computational work was subsequently built. His foundational work in "Causality: Models, Reasoning, and Inference" \cite{pearl2000causality} was as influential as his foundational critique in "Theoretical Impediments to Machine Learning with Seven Sparks from the Causal Revolution" \cite{pearl2018theoretical}. In particular, his contribution to the algorithmization of counterfactuals has proven instrumental for the development of contextual counterfactuals.

\textbf{Bareinboim's Transportability of Causal Effects}

Bareinboim's calculus of transportability \cite{bareinboim2012transportability} and his subsequent work on data fusion formalize the very problem of contextual variance that the EPP's explicit context is designed to manage at a computational level. Where Bareinboim provides the definitive logical framework for reasoning about moving causality between discrete contexts, the EPP provides the computational primitive, the dynamic, queryable, and multi-modal Context, to operationalize this reasoning as part of the EPP.

\textbf{Forbus's Defeasible Deontic Calculus}

Kenneth Forbus's work on formalizing deontic calculus \cite{olson2024DDIC} has proven invaluable to solve the complex topic of conflicting norms in the effect ethos. In practice, it is rarely possible to write conflict-free norms, and during the development of the effect ethos, a recurring theme was the acceptance of this reality. The subsequent search for a solution led to the adoption of Forbus's Defeasible Deontic Calculus as the primary means to resolve normative conflicts.

\textbf{Russel's Critique on Causality}

Bertrand Russell critique on causality  formulated in his 1912 essay ``On the Notion of Cause''\cite{RussellOnCause} directly lead to the realization that it's not necessarily causality that is central to Russels objection, but the underlying assumption of time asymmetry that is at odds with his observation that most successful theories in physics are based on time symmetry. Russel hinted at a profound truth because while physics routinely models dynamic change in complex systems, computational causality consistently struggles with capturing dynamic causality.  While there is truth to structural invariance, there is also the reality that dynamic systems emit different causal structures depending on dynamic change and that is where computational causality is fundamentally at odds with physics and to an extent with reality. From there, it became clear that for causality to handle dynamics requires a new foundation of causality itself.

\textbf{Hardy's Causaloid}

Lucian Hardy introduced the "causaloid,"\cite{HardyDynamicCausalStructure} a concept that encapsulates a spatial region and the causal connections within 
as a foundation for his work on finding a theory of Quantum Gravity. Critically, unlike all prior forms of causality, 
Hardy's causaloid is spacetime-agnostic because it folds cause and effect into one entity and thus removes the need for temporal order. His seminal work "Probability Theories with Dynamic Causal Structure" \cite{hardy2005probability} had a three-fold impact.
First, during the foundational work of lifting causality into geometric structures, his causaloid formalism has
proven instrumental in the formation of isomorphic recursive causal data structures. 
Second, his insight that the causaloid formalism puts deterministic and probabilistic structures on equal footing directly led to the multi-modal reasoning of the EPP. 
Third, his demonstration that fundamental differences of theoretical foundations are contained in a causaloid directly informed the representation of causal relations 
as a causal function, which then resulted in the single axiomatic formulation of the EPP.

\newpage

\textbf{Bornholdt's Uncertain Type}

Bornholdt's et al. contribution in "Uncertain< T >: A First-Order Type for Uncertain Data" \cite{bornholt2014uncertain} directly informed the unification of deterministic and probabilistic reasoning in the EPP. Instead of representing a value with a single number, the Uncertain type in the EPP represents a probability with a full probability distribution (e.g., Normal, Bernoulli) or even a complex computation graph that produces a distribution. More importantly, the EPP reasoning logic can seamlessly "lift" simpler Deterministic and Probabilistic effects into the Uncertain distribution, aggregate all distributions, and then infer a logical combination of all the input distributions without loss of information. The final output is an Uncertain<bool> that collapses rich uncertainties into a single value at the very last moment while preserving second-order properties such as the standard deviation or confidence level. As a result, decisions, and crucially, deontic reasoning become more robust under uncertainty.


\subsection{Contribution}

The Effect Propagation Process contribution spans three distinctive areas. First and foremost a single axiomatic foundation  formally derives the entire EPP from first principles. Second, the EPP establishes dynamic and emergent causal structures and an accompanying computable ethics framework for verifiable safety of dynamic causal systems. Third, a reference implementation in Rust demonstrates dynamic causality.   


\subsubsection{A Single Axiomatic Foundation of Dynamic Causality}

The EPP's single axiom, causality as a spacetime-agnostic functional dependency generalizes causality so that the classical definition becomes a special case of it. Section \ref{sec:philosophy} establishes the underlying philosophical foundation of the EPP. The details of the single axiomatic foundation are established in section \ref{sec:epp_definition}, the causal invariant of mechanism autonomy that supplies causal content beyond the monadic form is established in section \ref{sub:epp_causal_invariant}, the implications of the generalization are discussed throughout section \ref{sec:metaphysics} and \ref{sec:teleology}, and the formalization substantiates the EPP in section \ref{sec:formalization} further. The generalization is further substantiated in section \ref{sec:epp_subsumption_classical} where five established methods of classical causality are expressed through the EPP  and in Appendix A with a formal proof that the EPP subsumes the SCM. The SCM was chosen for the proof because of its historical and foundational significance to the field of computational causality. The reframing of causality as a single axiomatic functional dependency establishes explicit context and the causaloid as first-class structures. On important conjecture that follows from the axiom of functional dependency is that the field of function theory becomes applicable to causality as derived in section \ref{sub:epp_general_def_causality} and further discussed in section \ref{sec:future_work_functional_causality}. The impact of the new foundation results in a single, coherent language to model advanced dynamic causality in physics, software, finance, and system biology with equal rigor.

\subsubsection{A Formalism for Dynamic and Emergent Causal Structures}


The EPP is designed form the ground up to model meta-causality, or the causality of causal change. The EPP defines an operational Generative Function is a function that can modify the causal model, its context, or both dynamically to describe causal emergence. Section \ref{sec:metaphysics} establishes the foundation of causal emergence, the higher-order implications are discussed in section \ref{sec:teleology}, and section \ref{sec:formalization_epp} establishes the formalization. The impact of causal emergence enables systems that can generate novel causal strategies in response to unforeseen circumstances. When safeguarded by the Effect Ethos, this provides a path to creating robust and resilient systems that can safely navigate a dynamically changing world.

\subsubsection{A Computable Ethics Framework for Verifiable Safety}

Traditional causal models are overwhelmingly focused on passive inference and prediction. They provide a powerful way to answer "what if" questions but lack a formal, verifiable, and safe bridge from that knowledge to taking action. Safety, ethics, and alignment are treated as post-hoc problems to be solved with testing and ad-hoc rules.

The EPP is the first computational causality framework to treat verifiable safety as first-class formal primitives directly integrated into the foundation itself. The motivation for the Effect Ethos is rooted in safeguarding causal emergence, and the EPP achieves this with two mechanisms:

\begin{itemize}
\item The Causal State Machine (CSM):  Links a causal inference to a deterministic action.
\item The Effect Ethos: The Effect Ethos allows a system to formally verify that a proposed action is compliant with its mission and safety objectives before it is executed.
\end{itemize}

The impact of the Effect Ethos results in a feasible roadmap for building trustworthy, high-stakes autonomous systems that adhere to contextual rules of engagement encoded in an immutable ethos. The effect ethos inverts the entire verification and validation model by shifting the safety objectives of a system into the pre-design stage to convert existing rules of engagement into a testable effect ethos, which enables design-time safety and provable alignment.

\subsubsection{A Reference Implementation in Rust}

The presented Effect Propagation Process has been fully implemented in the open\-source DeepCausality\footnote{\url{www.deepcausality.com}} project hosted at the Linux Foundation for Data \& AI. The implementation comprises a set of crates that together realize the framework and contribute independently usable artifacts to the Rust ecosystem (described below). The five established methods of classical causality expressed through the EPP in section \ref{sec:epp_subsumption_classical} are demonstrated as code examples\footnote{\url{https://github.com/deepcausality-rs/deep_causality/tree/main/examples}} in the DeepCausality project. At the time of writing, the DeepCausality mono-repository has reached a total of 200,000 lines of Rust code across twenty crates with a sustained three-month average test coverage\footnote{\url{https://app.codecov.io/gh/deepcausality-rs/deep_causality}} of ninety-five percent.


The reference implementation comprises a layered scientific-computing stack with the EPP framework sitting on top. Each layer is an independently usable artifact within the Rust ecosystem, and the layers compose through the higher-kinded type machinery established in the foundation layer. The stack consists of a foundation layer, three mathematical layers built on the foundation, a physics layer built on the mathematical layers, a causal layer that hosts the EPP framework, and a causal-discovery layer at the top, with a horizontal metric crate providing a single source of truth for sign conventions across all layers. The collection is not exhaustive in the sense of covering every operation a mature standard library would provide, but every implemented operation is grounded in the corresponding mathematical structure and tested against established reference results.
 
\textit{Foundation layer} 

Three crates establish the foundation the DeepCausality project is built upon. The numerical foundation library (\texttt{deep\_causality\_num}) provides a complete trait hierarchy for abstract algebra, encoding the categorical structure of mathematical objects as a Rust type system. The hierarchy ranges from foundational structures (Magma, Semigroup, Monoid) through groups (Group, AbelianGroup) and rings (Ring, AssociativeRing, CommutativeRing, EuclideanDomain) to fields (Field, RealField, ComplexField) and module-and-algebra structures (Module, Algebra, AssociativeAlgebra, DivisionAlgebra). The hierarchy is wired so that types implementing each level satisfy the corresponding mathematical laws as compile-time assertion. The crate also provides Float106, a 106-bit high-precision floating-point type with approximately thirty-two decimal digits of precision, which is comparable to IEEE binary128 but available now on stable Rust and several times faster on most hardware. Full implementations of Complex, Quaternion, and Octonion as division algebras are included. The library is macro-free, unsafe-free, and dependency-free, and supports no-std environments. 

The higher-kinded type (\texttt{deep\_causality\_haft}) crate provides a working implementation of arity-5 higher-kinded types in Rust through type-level encoding and the witness pattern. The crate supplies the Effect, Functor, Applicative, Monad, and CoMonad trait families that all higher layers build on. Rust does not provide higher-kinded types natively, and the haft crate is what allows the mathematical and causal crates to compose monadically by implementing the HKT traits.

The metric crate (\texttt{deep\_causality\_metric}) provides a single source of truth for metric signatures across the stack. It defines the East Coast convention (used in general relativity), the West Coast convention (used in particle physics), and the broader Clifford signature space $Cl(p, q, r)$ that the multivector and topology crates build on. The horizontal placement of this crate is what allows a particle-physics calculation in one layer and a general-relativity calculation in another layer to share compatible sign conventions without manual coordination, eliminating an entire class of consistency bugs.
 
\textit{Mathematical layers} 

Three crates implement distinct mathematical structures over the foundation, each with higher-kinded type instances that allow composition with the others and with the EPP framework. The tensor crate (\texttt{deep\_causality\_tensor}) provides multi-dimensional arrays with stride-based memory layout, broadcasting for element-wise operations, Einstein summation for matrix products and contractions, and Functor, Applicative, Monad, and CoMonad instances for functional composition. 

The multivector crate (\texttt{deep\_causality\_multivector}) provides Clifford algebras over the dynamic signature space, with pre-configured constructors for Pauli algebra, spacetime algebra, conformal geometric algebra, projective geometric algebra in three dimensions, the Dixon algebra used in Standard Model particle physics, and the Grand Unified Algebra hosting the full $\mathfrak{spin}(10)$ gauge symmetry. The crate implements Functor, Applicative, and Monad instances where the bind operation realizes the tensor product of algebras, allowing dimension-changing compositions to be expressed monadically. 

The topology crate (\texttt{deep\_causality\_topology}) provides graphs, hypergraphs, simplicial complexes, manifolds, and point clouds, together with differential geometry operators (exterior derivative, Hodge star, codifferential, Hodge-Laplacian) and a lattice gauge field framework supporting U(1), SU(2), SU(3), and Lorentz gauge groups. The lattice gauge implementation is verified against twenty-four reference results from Creutz's \emph{Quarks, Gluons and Lattices}\cite{creutz1983quarks}. The crate implements Functor and BoundedComonad instances, which makes graph convolutions and cellular automata expressible through the comonadic extend operation.
 
\textit{Physics layer} 

The physics library (\texttt{deep\_causality\_physics}) builds on the four layers below it, with kernels organized into thematic domains spanning astrophysics, condensed matter, classical dynamics, electromagnetism, fluid dynamics, materials science, magnetohydrodynamics, nuclear physics, photonics, quantum mechanics, relativity, thermodynamics, and wave mechanics with all using type-safe units. The library separates two levels of usage. First, physics kernels are pure stateless functions for isolated equations such as the Schwarzschild radius, the Lorentz force, the Klein-Gordon equation, or the Lund string fragmentation. Second, each kernel has a corresponding wrapper that wraps the return type into a propagating effect for usage and composition in the effect propagation process. Theories are full physical frameworks (general relativity, electromagnetism, weak force, electroweak unification) implemented as gauge theories over the topology crate's gauge field abstraction. The theory layer is what allows different physical theories to share the same field-strength computation, the same parallel transport, and the same Bianchi-identity machinery, with the gauge group as the only differentiator. The physics library is also generic over the floating-point type, which means a simulation can be run at single, double, or Float106 precision by changing a single type alias.
 
\textit{Causal layer} 

Four crates constitute the EPP framework. The core crate (\texttt{deep\_causality\_core}) provides the foundational types of the EPP framework: PropagatingEffect for stateless causal effects, PropagatingProcess for stateful effects with state and configuration channels, the CausalMonad with its bind operation, and the Causal Effect System runtime that coordinates causal models. The same crate also provides the ControlFlowBuilder, a parallel capability for constructing compile-time-verified static execution graphs in which the Rust type system enforces that nodes can only be connected if their input and output types match. With the optional \texttt{strict-zst} feature, the ControlFlowBuilder enforces zero-sized types throughout the graph, which eliminates dynamic dispatch and heap allocation in the execution path and produces graphs amenable to worst-case execution time analysis and formal verification, suitable for safety-critical and hard real-time deployment. 

The main crate (\texttt{deep\_causality}) builds on the core and on the wider stack to provide the Causaloid in its three isomorphic forms (Singleton, Collection, Graph), the Context with its Euclidean and non-Euclidean representations, the Causal State Machine, the Teloid and Effect Ethos for deontic governance, and the integration of all of these into the EPP framework. UltraGraph (\texttt{ultragraph}) is a high-performance two-phase hypergraph data structure that supports the CausaloidGraph and the Context Hypergraph and is independently usable wherever performance-constrained hypergraph composition is needed. The Uncertain Type (\texttt{deep\_causality\_uncertain}) is a Rust implementation of Bornholt's Uncertain<T> as a first-order type for uncertain data\cite{bornholt2014uncertain}.

The composition through higher-kinded types means that a tensor of multivectors over a topology defined as a lattice gauge field, with each layer respecting the algebraic laws of the underlying types and all sharing the metric crate's sign conventions, is expressible as a single type and operates monadically across the stack.
 
 \textit{Causal discovery layer} 
 
The Causal Discovery Language (\texttt{deep\_causality\_discovery}) is a domain-specific language that bridges from raw observational data to executable causal models. The pipeline uses Rust's typestate pattern to encode the sequence of stages (data loading, cleaning, feature selection, causal discovery, analysis, finalization) so that misuse is rejected at compile time. The discovery language closes the loop from observational data to causal model that is left open in most causal frameworks.
 
  
\subsubsection{Code Examples}

The examples chapter (Chapter\ref{sec:examples}) presents three examples that demonstrate different capabilities of the framework at progressively more advanced levels. The first, the event horizon probe, demonstrates regime change in a setting where the causal structure itself evolves as the system crosses a physical threshold (§\ref{sec:example_event_horizon}). The second, the flight envelope monitor, demonstrates the compositional capability described above by running a single pipeline that interleaves Causaloid evaluations with CausalMonad bind steps, with state and audit log accumulating across both kinds of containers (§\ref{sec:example_avionics}). The third demonstrates multi-stage multi-physics composition through a chronogravimetric application that recovers the geocentric gravitational constant from satellite-clock time-dilation data (§\ref{sec:example_chronogravimetry}). Chronogravimetry inherits its core relativistic-redshift methods from the chronometric-geodesy literature initiated by Bjerhammar (1975) and Vermeer (1983). 



\subsection{Structure}

The presented monograph serves multiple audiences.

\textbf{Primary audience:} Researchers in computational causality, causal inference, and AI safety who are reaching the limits of classical causal frameworks and need a foundation for dynamic and emergent causal structures. For these readers a natural path is to begin with the philosophy of causality (Chapter \ref{sec:philosophy}) and the overview of the EPP (Chapter \ref{sec:epp}), and then delve into its philosophical foundation: the Metaphysics (Chapter \ref{sec:metaphysics}), the Epistemology (Chapter \ref{sec:epp_epistemology}), and the Teleology (Chapter \ref{sec:teleology}).

\textbf{Secondary audience:} Practitioners in avionics, materials science, theoretical physics, and other safety-critical or precision-critical engineering disciplines who need a verifiable computational substrate for causal reasoning under dynamic conditions. The reference implementation is what brings them in. These readers may start with the high-level EPP framework (Chapter \ref{sec:epp}) and then proceed to the DeepCausality Implementation chapter (Chapter \ref{sec:implementation}), which is organized in layered fashion: foundation, mathematical layers, physics, causal framework, and causal discovery. Engineering readers may locate the layer most relevant to their domain and read outward from there. The Metaphysics chapter (Chapter \ref{sec:metaphysics}) provides details about the orthogonal design used throughout the implementation.


\textbf{Tertiary audience:} Philosophers of science working on causation, process metaphysics, and the philosophical foundations of computation, who will engage with the metaphysics, epistemology, and teleology chapters as substantive philosophy.


\newpage